\chapter{Who Gets to Talk to the Plant?}

Picture a chemical plant with one data server sitting next to it. Two programs
want to know what the plant is doing. One draws screens for the operators. The
other keeps history for the engineers. Each opens its own connection, asks for
the same numbers a couple of times per second, and stores what it gets in its
own place. Neither knows the other exists.

\begin{figure}[htbp]
\centering
\begin{tikzpicture}
  \node[hf/source, minimum width=36mm] (srv) {Plant data server};
  \node[hf/reader, minimum width=32mm, below left=16mm and 4mm of srv] (hmi) {Operator screens};
  \node[hf/reader, minimum width=32mm, below right=16mm and 4mm of srv] (his) {History program};
  \node[hf/store, below=9mm of hmi] (s1) {own copy};
  \node[hf/store, below=9mm of his] (s2) {own copy};

  \draw[hf/ask] (hmi.north) -- node[hf/note, left=2pt] {asks for everything,\\twice a second} (srv.south west);
  \draw[hf/ask] (his.north) -- node[hf/note, right=2pt] {asks for everything,\\twice a second} (srv.south east);
  \draw[hf/save] (hmi) -- (s1);
  \draw[hf/save] (his) -- (s2);
  \draw[hf/pain, dashed] (hmi.east) -- node[hf/note, below=2pt, text=hfAlert] {unaware of\\each other} (his.west);
\end{tikzpicture}
\caption{The situation: two readers, one source, two copies of the truth.}
\end{figure}

Nothing here is broken. Both programs work. But there are now two copies of the
truth, two loads on the server, and two different answers to ``what happened at
three in the morning?''. Before adding a third program, it is worth asking a
question that industry has been answering, and re-answering, for more than fifty
years: \emph{who gets to talk to the plant?}

This chapter follows that question from relay cabinets to message brokers. Each
answer was forced by a specific pain, and each one fixed something the previous
answer got wrong.

\begin{figure}[htbp]
\centering
\begin{tikzpicture}[x=0.23cm]
  \draw[black!40, line width=1pt, -{Stealth[length=2mm]}] (0,0) -- (61,0);
  \foreach \y in {1970,1980,...,2020} {
    \draw[black!40] (\y-1965,0.08) -- (\y-1965,-0.08);
  }
  % year / height (+ above, - below) / color / label
  \foreach \yr/\h/\c/\t in {
      1968/0.5/hfSource/{\textbf{1968}\\first PLC},
      1979/-0.5/hfSource/{\textbf{1979}\\Modbus},
      1980/1.5/hfStore/{\textbf{1980}\\OSIsoft and PI},
      1990/0.5/hfSource/{\textbf{1990}\\Tennessee Eastman},
      1996/-0.5/hfBus/{\textbf{1996}\\OPC DA},
      1999/1.5/hfBus/{\textbf{1999}\\MQTT},
      2000/-1.5/hfStore/{\textbf{2000}\\ISA-95},
      2008/0.5/hfBus/{\textbf{2008}\\OPC UA},
      2014/-0.5/hfBus/{\textbf{2014}\\MQTT at OASIS},
      2021/0.5/hfStore/{\textbf{2021}\\AVEVA buys OSIsoft}} {
    \fill[\c] (\yr-1965,0) circle (2.2pt);
    \draw[\c!70] (\yr-1965,0) -- (\yr-1965,\h);
    \pgfmathsetmacro\up{\h>0 ? 1 : 0}
    \ifdim\up pt>0pt
      \node[font=\sffamily\scriptsize, text width=19mm, align=center, anchor=south, inner sep=2pt] at (\yr-1965,\h) {\t};
    \else
      \node[font=\sffamily\scriptsize, text width=19mm, align=center, anchor=north, inner sep=2pt] at (\yr-1965,\h) {\t};
    \fi
  }
\end{tikzpicture}
\caption{The milestones in this chapter. Gray: devices and plants. Purple:
something in the middle. Green: memory.}
\end{figure}

\section{A plant that never existed}

The plant in this story is fictional, and its fictionality is a piece of history
in its own right.

\begin{history}{1990 \textperiodcentered{} Kingsport, Tennessee}
James J.\ Downs and Ernest F.\ Vogel, two engineers at Eastman Chemical Company,
presented a simulated process at an AIChE meeting. The paper came out in 1993 in
\emph{Computers \& Chemical Engineering} under the title ``A plant-wide
industrial process control problem''.

The problem they were solving was social rather than chemical. Academics who
studied control had nothing realistic to test on, and companies would never
publish the model of a real plant, because that model \emph{is} the company's
know-how. So Eastman built a plausible plant with the chemistry disguised:
components are called A to H, and the reactions are just ``gas-liquid,
exothermic''.
\end{history}

The Tennessee Eastman Process has a reactor, a condenser, a vapor-liquid
separator, a recycle compressor and a stripper. It exposes 41 measured variables
and 12 manipulated ones, plus 20 disturbances you can switch on to make life
harder. It is open-loop unstable: leave it without control and it drifts until
its safety limits shut it down. That made it the standard benchmark for
plant-wide control, fault detection and, decades later, machine learning on
process data.

There's one detail worth keeping for later. Not every measurement moves at the
same pace. Most are continuous, but the composition analyzers report only every
0.1 or 0.25 hours of plant time, after a dead time. A program that asks for
``everything'' twice a second will mostly hear the same answers repeated, and
that will matter later in the chapter.

\section{The oldest way to know is to ask}

\begin{history}{1968--1979 \textperiodcentered{} Detroit and Massachusetts}
General Motors' Hydra-Matic division was tired of rewiring relay cabinets every
time a production line changed. In 1968 it wrote a specification for an
electronic replacement. Dick Morley's team at Bedford Associates answered with
the Modicon~084, and the programmable logic controller (PLC) was born.

PLCs soon needed to talk to something above them. In 1979, Modicon published
\textbf{Modbus}: on a serial line, one master asks and the slaves answer. In
2020, the Modbus Organization renamed the roles to client and server, but the
rule is the same.
\end{history}

Under Modbus, a slave never speaks unless spoken to. If you want to know a
temperature, you ask for register 40001, and if you want to know it again, you
ask again.

\begin{figure}[htbp]
\centering
\begin{tikzpicture}[x=1cm, y=-0.72cm]
  \node[hf/reader, minimum width=22mm] (m)  at (0,0)   {Master};
  \node[hf/source, minimum width=22mm] (s1) at (4.8,0) {Slave 1};
  \node[hf/source, minimum width=22mm] (s2) at (9.6,0) {Slave 2};
  \foreach \x in {0,4.8,9.6} \draw[hf/lifeline] (\x,0.7) -- (\x,8.6);

  \draw[hf/ask]    (0,1.4) -- node[hf/note, above] {read 40001} (4.8,1.4);
  \draw[hf/answer] (4.8,2.2) -- node[hf/note, above] {21.4} (0,2.2);
  \draw[hf/ask]    (0,3.2) -- node[hf/note, above, pos=0.3] {read 40001} (9.6,3.2);
  \draw[hf/answer] (9.6,4.0) -- node[hf/note, above, pos=0.7] {57.0} (0,4.0);
  \draw[hf/ask]    (0,5.2) -- node[hf/note, above] {read 40001} (4.8,5.2);
  \draw[hf/answer] (4.8,6.0) -- node[hf/note, above] {21.4 \emph{again}} (0,6.0);
  \draw[hf/ask]    (0,7.0) -- node[hf/note, above, pos=0.3] {read 40001} (9.6,7.0);
  \draw[hf/answer] (9.6,7.8) -- node[hf/note, above, pos=0.7] {57.0 \emph{again}} (0,7.8);

  \node[hf/note, anchor=west] at (10.0,4.5) {time\\$\downarrow$};
\end{tikzpicture}
\caption{Polling on a Modbus line. The master asks in turn. Every answer costs a
round trip, even when nothing changed.}
\end{figure}

\begin{concept}{Polling}
The reader asks the source on a schedule, and the source answers and forgets.
The source is passive. It doesn't know who is interested or how often.

\textbf{What it fixes:} order on a shared wire. With one party starting every
conversation, nothing collides.

\textbf{What it costs:} work grows with readers $\times$ frequency, whether or
not anything changed. Fast changes between two questions are lost, and slow
values get read over and over.
\end{concept}

Polling made sense at the time. On a shared serial wire with no collision
handling, it is what keeps the line orderly. It is the native shape of the field
bus.

It also leaves a mark that's still with us: the source is passive. If a second
program wants the same temperature, it has to ask too. On a Modbus line it
can't, because there is only one master, so the second program has to go
through the first.

\section{Every screen with its own drivers}

During the 1980s, operator screens moved from dedicated consoles to ordinary
PCs. HMI and SCADA packages appeared, and each of them had to read from PLCs
made by many vendors, each with its own protocol.

\begin{concept}{The N$\times$M problem}
When N programs need to talk to M kinds of device and each pair needs its own
code, you write N$\times$M pieces of code. If you put a standard interface in
the middle, each side implements it once, and the count drops to N$+$M.
\end{concept}

Every HMI vendor wrote a driver for every device family. Ten HMIs and twenty
device families meant two hundred drivers, written separately and often badly,
and all of them polling.

\begin{figure}[htbp]
\centering
\begin{tikzpicture}
  \begin{scope}
    \foreach \i/\n in {0/A,1/B,2/C} {
      \node[hf/reader, minimum width=15mm] (h\i) at (0,-1.2*\i) {HMI \n};
      \node[hf/source, minimum width=15mm] (d\i) at (3.6,-1.2*\i) {PLC \the\numexpr\i+1\relax};
    }
    \foreach \i in {0,1,2} \foreach \j in {0,1,2}
      \draw[hf/pain, draw=hfAlert!70] (h\i.east) -- (d\j.west);
    \node[hf/note, text=hfAlert] at (1.8,-3.3) {N$\times$M $=$ 9 drivers};
  \end{scope}
  \begin{scope}[xshift=7.2cm]
    \foreach \i/\n in {0/A,1/B,2/C} {
      \node[hf/reader, minimum width=15mm] (k\i) at (0,-1.2*\i) {HMI \n};
      \node[hf/source, minimum width=15mm] (e\i) at (4.4,-1.2*\i) {PLC \the\numexpr\i+1\relax};
    }
    \node[hf/bus, minimum height=34mm, minimum width=10mm] (opc) at (2.2,-1.2) {\rotatebox{90}{OPC interface}};
    \foreach \i in {0,1,2} {
      \draw[hf/flow, draw=hfBus] (k\i.east) -- (k\i.east -| opc.west);
      \draw[hf/flow, draw=hfBus] (e\i.west -| opc.east) -- (e\i.west);
    }
    \node[hf/note, text=hfBus] at (2.2,-3.3) {N$+$M $=$ 6 implementations};
  \end{scope}
\end{tikzpicture}
\caption{Left: every HMI talks to every device in its own way. Right: both sides
speak one standard.}
\end{figure}

\begin{history}{1995--1996 \textperiodcentered{} The OPC task force}
A group of automation vendors, among them Fisher-Rosemount, Intellution, Opto~22
and Rockwell Software, got together with Microsoft to fix the driver mess. They
released the \textbf{OPC Data Access} specification in 1996. At the time, OPC
stood for \emph{OLE for Process Control}, and it was built on Microsoft's COM and
DCOM.
\end{history}

COM and DCOM are names that only make sense next to the names that came before
them, so it's worth a short detour through Windows.

\begin{history}{1987--1996 \textperiodcentered{} Redmond, and one HMI company}
\textbf{1987, DDE.} Windows~2.0 brought \emph{Dynamic Data Exchange}, the first
way for two Windows programs to swap data. ``Dynamic'' meant the link stayed
alive: change a cell in Excel and the other program got the new value. An HMI
company, Wonderware, wrote \emph{NetDDE} to carry DDE across a network, and
Microsoft licensed it for Windows for Workgroups. Before OPC, this is how many
HMIs got their numbers.

\textbf{1990, OLE.} \emph{Object Linking and Embedding} is named after its two
tricks: \emph{embed} a copy of an Excel chart inside a Word document, or
\emph{link} to the file so the chart follows its changes. Double-click the chart
and Excel opens inside Word to edit it. Now Word had to \emph{call Excel's code},
not just receive a value.

\textbf{1993, COM.} For OLE~2, Microsoft pulled the generic machinery out and
gave it its own name, the \emph{Component Object Model}. \emph{Component}: a
binary piece that programs can use without having been compiled with it.
\emph{Object}: something with methods behind interfaces. \emph{Model}: a set of
rules, not a library.

\textbf{1996, DCOM.} First called \emph{Network OLE}, it shipped with Windows
NT~4.0 as \emph{Distributed COM}: the same calls, to an object on another
machine.

So when the vendors called their standard \emph{OLE for Process Control}, they
meant: use the Windows machinery for one program to call another, and agree on
the interfaces for reading plant values.
\end{history}

\begin{concept}{COM, the component object model}
An \emph{interface} is a contract: a named list of methods. Each interface is
identified by a \emph{GUID}, a 128-bit \emph{globally unique identifier}
(Microsoft's name for the UUIDs of the DCE world, more on DCE below). The
Windows registry maps each GUID to the DLL or program that implements it.

Every COM object implements one interface called \texttt{IUnknown}. It's named
that because, when you get hold of an object, you don't yet know what it is. You
ask it, through \texttt{QueryInterface}: ``do you speak \texttt{IOPCServer}?''.
If the answer is yes, you get that interface and call its methods.

\textbf{What it fixes:} programs from different companies, in different
languages, call each other without sharing source code.

\textbf{What it costs:} everything rests on the Windows registry and the Windows
binary format.
\end{concept}

DCOM adds one thing on top of COM: the object can live on another machine, and
the caller doesn't notice. That trick is older than Windows.

\begin{history}{1976--1994 \textperiodcentered{} The remote procedure call}
In 1976, James E.\ White's RFC~707 proposed that programs share resources over
a network as if they were calling procedures. In 1984, Andrew Birrell and Bruce
Nelson at Xerox PARC published ``Implementing Remote Procedure Calls'', which
showed how to make a call to another machine look like an ordinary local call.
Sun built its RPC under NFS. The Open Software Foundation's \emph{Distributed
Computing Environment} (DCE) standardized another one in the early 1990s, and
Microsoft adopted it as MSRPC, the layer under DCOM.

In 1994, Jim Waldo and colleagues at Sun Labs wrote ``A Note on Distributed
Computing''. They argued that pretending a remote call is local is a mistake,
because a remote call can be slow, fail halfway, or leave you not knowing
whether it ran.
\end{history}

\begin{figure}[htbp]
\centering
\begin{tikzpicture}
  \node[hf/reader, minimum width=24mm] (hmi) {HMI};
  \node[hf/base, draw=hfReader, fill=white, minimum width=24mm, below=10mm of hmi] (px) {proxy};
  \node[hf/base, draw=hfSource, fill=white, minimum width=24mm, right=62mm of px] (st) {stub};
  \node[hf/source, minimum width=24mm, above=10mm of st] (srv) {OPC server};

  \begin{scope}[on background layer]
    \node[draw=black!30, rounded corners=3pt, inner xsep=9mm, inner ysep=4mm, fit=(hmi)(px),
          label={[hf/note]below:machine A}] {};
    \node[draw=black!30, rounded corners=3pt, inner xsep=9mm, inner ysep=4mm, fit=(st)(srv),
          label={[hf/note]below:machine B}] {};
  \end{scope}

  \draw[hf/ask] (hmi) -- node[hf/note, left] {\texttt{Read(...)}} (px);
  \draw[hf/push] (px) -- node[hf/note, above] {RPC over the network} node[hf/note, below] {port 135, then a dynamic port} (st);
  \draw[hf/ask] (st) -- node[hf/note, right] {\texttt{Read(...)}} (srv);
\end{tikzpicture}
\caption{A DCOM call. The HMI calls a local proxy that looks exactly like the
server. The proxy packs the arguments and sends them, and a stub on the other
machine unpacks them and makes the real call.}
\end{figure}

\begin{concept}{Remote procedure call}
The caller calls a local \emph{proxy} with the same signature as the real
object. The proxy \emph{marshals} the arguments, packing them into bytes, and
sends them. On the other machine, a \emph{stub} unpacks them, calls the real
object, and sends the result back the same way.

\textbf{What it fixes:} the network code disappears from the program. Calling a
server on another machine looks like calling a function.

\textbf{What it costs:} the network doesn't disappear. The call can hang, fail
halfway or time out, and the code that looked local has to handle all of it.
\end{concept}

Back to OPC. The device vendor writes \emph{one} OPC server, and any OPC client
can read from it. OPC DA also brought something Modbus didn't have, the subscription. A client
could create a \emph{group} of items with an update rate, and the server would
call the client back when values changed. The server still polled the device
underneath, but it did so once, on behalf of everyone.

The price was the foundation it stood on. COM and DCOM meant Windows only. DCOM
across a network meant fighting firewalls: it first asks port 135 where the
object lives, then talks on a port chosen at run time. Permissions lived in a
tool called \texttt{dcomcnfg}, per machine and per user, and a wrong combination
failed silently. A whole profession of ``DCOM configuration'' grew up around it.

Microsoft moved on to .NET in 2002, and DCOM became legacy. Between 2021 and
2023, Microsoft rolled out a security hardening of DCOM, and when it was
enforced, OPC Classic connections stopped working in plants that hadn't
prepared. By then the industry had long since started building the successor.

\section{OPC UA: the same idea without Windows}

\begin{history}{2008 \textperiodcentered{} OPC Unified Architecture}
Work on a successor started in the early 2000s. In 2008, the OPC Foundation
released \textbf{OPC UA}, later standardized as IEC~62541. In 2011, the
foundation kept the letters but changed what they stood for: \emph{Open Platform
Communications}.
\end{history}

UA kept what worked in OPC and rebuilt the rest. It has its own binary protocol
over TCP, so it can run on Linux, inside a PLC, or in a container. It has
security built in. It also has a real information model: the server exposes
typed nodes in folders, not just a flat list of tags.

\begin{concept}{Subscription}
Instead of asking for values, the reader tells the source what it cares about
and how often, and the source reports changes. In OPC UA, a client creates a
\emph{subscription} with a \emph{publishing interval} and adds \emph{monitored
items} to it. Each item has a \emph{sampling interval} and, optionally, a
\emph{deadband}: a change smaller than that is not worth reporting.

\textbf{What it fixes:} the source does the sampling once and sends only what
changed, so a slow value costs almost nothing.

\textbf{What it costs:} the source now keeps state per client, and that state
has limits.
\end{concept}

How UA delivers those changes is a neat trick. UA is a request/response
protocol, so the server can't really speak first. Instead, the client sends
\texttt{Publish} requests \emph{in advance}, and the server holds them until it
has notifications to send. When it answers one, the client sends another.

\begin{figure}[htbp]
\centering
\begin{tikzpicture}[x=1cm, y=-0.75cm]
  \node[hf/reader, minimum width=24mm] at (0,0) {Client};
  \node[hf/source, minimum width=24mm] at (7,0) {Server};
  \draw[hf/lifeline] (0,0.7) -- (0,7.4);
  \draw[hf/lifeline] (7,0.7) -- (7,7.4);

  % held requests
  \fill[hfBus!25] (6.88,1.6) rectangle (7.0,3.9);
  \fill[hfBus!45] (7.0,2.6) rectangle (7.12,6.6);

  \draw[hf/ask] (0,1.2) -- node[hf/note, above, sloped] {Publish \#1} (7,1.6);
  \draw[hf/ask] (0,2.2) -- node[hf/note, above, sloped] {Publish \#2} (7,2.6);
  \node[hf/note, anchor=west] at (7.3,2.1) {both held:\\nothing to say yet};

  \node[hf/note, anchor=west, text=hfBus] at (7.3,3.6) {a value changes};
  \draw[hf/push] (7,3.9) -- node[hf/note, above, sloped] {answer \#1: the change} (0,4.3);
  \draw[hf/ask] (0,4.9) -- node[hf/note, above, sloped] {Publish \#3} (7,5.3);
  \node[hf/note, anchor=west, text=hfBus] at (7.3,6.3) {another change};
  \draw[hf/push] (7,6.6) -- node[hf/note, above, sloped] {answer \#2} (0,7.0);
\end{tikzpicture}
\caption{Push built out of requests. The client keeps a few \texttt{Publish}
requests waiting at the server, so the server always has one to answer when
something changes.}
\end{figure}

That is push built out of polling, the same move the web would make with long
polling years later.

So a UA server can serve many clients efficiently, but only if they subscribe.
If every client simply reads the whole address space on a timer, the server does
the full work once per client per tick, and the analyzer that changes every six
minutes of plant time gets read a few thousand times in between. Embedded UA
servers also cap the number of sessions and subscriptions, sometimes to only a
handful. Every new reader that goes straight to the source uses up capacity of the
server, and nobody keeps track of how much is left.

\section{The plant's memory}

Control systems were built to control, not to remember. A distributed control
system kept a short trend, perhaps hours, perhaps a few days. But engineers
wanted months, to compare this summer with last summer, to find the slow drift
before the failure, or to answer an auditor.

\begin{history}{1980 \textperiodcentered{} OSIsoft}
J.\ Patrick Kennedy founded Oil Systems, Inc., the company that became
\textbf{OSIsoft}, and built the \textbf{PI System}, short for \emph{Plant
Information}. PI defined the product category: the \emph{process historian}, a
time-series database built for plants before ``time-series database'' was a
phrase. In 2021, AVEVA bought OSIsoft.
\end{history}

A historian has to keep years of data from thousands of signals, and that runs
into a wall of arithmetic. Take a plant with 100{,}000 tags, each read once per
second:
\[
  100{,}000 \times 86{,}400 \times 365 \approx 3 \text{ trillion values a year.}
\]
At about 16 bytes each (timestamp, value, quality), that's around 50 terabytes a
year. When PI was born, a hard disk held tens of megabytes. Disks are cheap now,
but reading 50 terabytes to draw a one-month trend is still slow.

And almost all of it is repetition. A tank temperature sits at 80.0\,\textdegree
C for hours. Storing ``80.0'' 3{,}600 times an hour teaches nobody anything.
Storing a value that didn't change adds nothing. Historians act on that at two
points: before a value travels, and before it's written.

\begin{concept}{Exception reporting}
The collector sends a value only when it moves more than a \emph{deadband} away
from the last value it sent. With a deadband of 0.5\,\textdegree C, a tank
wandering between 79.8 and 80.2 sends nothing at all.

\textbf{What it fixes:} the network and the server stop carrying repetitions.

\textbf{What it costs:} changes smaller than the deadband are gone for good. So
the deadband is set per tag, close to the instrument's accuracy: a change
smaller than the sensor's own error is noise, not information.
\end{concept}

\begin{concept}{Compression}
Of the values that do arrive, the archive keeps only the ones needed to redraw
the curve with an error below a tolerance $E$. If a temperature climbs in a
straight ramp from 80 to 90 over an hour, two points are enough, the start and
the end, because a ruler redraws the middle. It's like drawing a road on a map:
you only mark the places where it bends.

The best-known algorithm is \emph{swinging door trending}, which E.~H.\ Bristol
at Foxboro patented in the 1980s and presented in 1990. At each new point it
asks one question: starting from the last stored point, can a single straight
line still pass within $E$ of every point since? Two ``doors'', hinged at $\pm
E$ above and below the stored point, keep that answer cheap, so the algorithm
never has to recheck old points. Each new point can only close the doors. When
they pass parallel, the answer is no.

\textbf{What it fixes:} the archive keeps the shape of the curve, not the
repetitions.

\textbf{What it costs:} the stored curve is exact only to within $E$. Between
two stored points, all you know is the straight line.
\end{concept}

\begin{figure}[htbp]
\centering
\newcommand\swaxes{%
  \draw[black!40, -{Stealth[length=1.6mm]}] (-0.15,0) -- (4.6,0);
  \draw[black!40, -{Stealth[length=1.6mm]}] (0,-0.1) -- (0,3.9);
}
\newcommand\swpt[2]{\draw[hfStore, fill=white, line width=0.8pt] (#1,#2) circle (2.1pt);}
\newcommand\swkeep[2]{\fill[hfStore] (#1,#2) circle (2.4pt);}
\begin{tikzpicture}[x=0.9cm, y=0.75cm]
  % 1. points arrive and still fit
  \begin{scope}
    \swaxes
    \node[hf/note, anchor=south west] at (0,3.9) {value};
    \node[hf/note, anchor=north east] at (4.6,0) {time};
    \fill[hfStore!12] (0,1.5) -- (3,1.8) -- (0,0.5) -- cycle;
    \draw[hfStore, dashed] (0,1.5) -- (3,1.8);
    \draw[hfStore, dashed] (0,0.5) -- (3,1.8);
    \node[hf/note, anchor=east] at (0,1.5) {$+E$};
    \node[hf/note, anchor=east] at (0,0.5) {$-E$};
    \draw[hfStore, densely dotted, line width=0.8pt] (0,1) -- (3,1.8);
    \swkeep{0}{1} \swpt{1}{1.3} \swpt{2}{1.5} \swpt{3}{1.8}
    \node[hf/note] at (2.1,-1.3) {\textbf{1.}\ Points arrive. One\\straight line still fits them all.};
  \end{scope}
  % 2. one point breaks the corridor
  \begin{scope}[xshift=5.6cm]
    \swaxes
    \fill[hfStore!12] (0,1.5) -- (3,1.8) -- (0,0.5) -- cycle;
    \draw[hfStore, dashed] (0,1.5) -- (3,1.8);
    \draw[hfStore, dashed] (0,0.5) -- (3,1.8);
    \draw[hfAlert, dashed, line width=0.8pt] (0,1) -- (4,3.4);
    \draw[hfAlert, line width=0.6pt, {Bar[width=1.6mm]}-{Bar[width=1.6mm]}] (3.18,1.8) -- (3.18,2.8);
    \node[hf/note, text=hfAlert, anchor=west] at (3.22,2.3) {$>E$};
    \swkeep{0}{1} \swpt{1}{1.3} \swpt{2}{1.5} \swpt{3}{1.8}
    \fill[hfAlert] (4,3.4) circle (2.4pt);
    \node[hf/note] at (2.1,-1.3) {\textbf{2.}\ A line to the red point\\misses the others by more than $E$.};
  \end{scope}
  % 3. what gets written
  \begin{scope}[xshift=11.2cm]
    \swaxes
    \draw[hfStore, line width=1pt] (0,1) -- (3,1.8);
    \draw[hfStore, densely dotted, line width=0.8pt] (3,1.8) -- (4,3.4);
    \draw[black!25, line width=0.6pt] (1,1.3) circle (2.1pt) (2,1.5) circle (2.1pt);
    \swkeep{0}{1} \swkeep{3}{1.8}
    \draw[hfStore, fill=white, line width=0.8pt] (4,3.4) circle (2.1pt);
    \node[hf/note, anchor=north west] at (3.05,1.7) {stored,\\new anchor};
    \node[hf/note] at (2.1,-1.3) {\textbf{3.}\ Two points written\\instead of four.};
  \end{scope}
\end{tikzpicture}
\caption{Swinging door trending in three steps. Filled points are stored, hollow
points are waiting, and gray points were dropped because the straight line
between the stored ones redraws them within $E$.}
\end{figure}

We've met this idea before: the deadband on an OPC UA monitored item is
exception reporting moved into the server, and MQTT will bring it back once
more.

Now look at a program that asks for everything twice a second and stores every
answer. It does the exact opposite of what fifty years of historians learned. It
keeps precisely the repetitions. An analyzer that changes every six minutes
becomes seven hundred identical rows between one real value and the next, which
is useless.

The historian also changed who talks to the plant. Once it existed, it was the
one program that already read everything, kept it, and knew how to answer
questions about the past. Other programs started asking the historian instead of
the plant. Security pushed in the same direction.

\begin{history}{Early 1990s \textperiodcentered{} Purdue University}
Theodore J.\ Williams and colleagues at Purdue published the \emph{Purdue
Enterprise Reference Architecture}. Its layered model, from field devices at
level~0 up to business systems at level~4, became the backbone of the ISA-95
standard (first part in 2000) and of industrial network segmentation.
\end{history}

The rule that came out of Purdue is simple: business networks don't talk to the
control network. The historian sits between them, and it's common to replicate
it into a DMZ so that the office reads a copy and never touches the plant.

\begin{figure}[htbp]
\centering
\begin{tikzpicture}[y=1.4cm]
  \foreach \y/\lvl/\desc/\shade in {
      5/Level 4/business/6,
      4/DMZ/between networks/14,
      3/Level 3/operations/6,
      2/Level 2/supervisory/6,
      1/Level 1/control/6,
      0/Level 0/process/6} {
    \fill[black!\shade] (0,\y-0.47) rectangle (11,\y+0.47);
    \node[font=\sffamily\small, anchor=west] at (0.2,\y) {\textbf{\lvl}\enspace{\color{black!60}\desc}};
  }
  \node[hf/reader, minimum height=7mm, minimum width=30mm] (off) at (6.6,5) {office apps};
  \node[hf/store, minimum height=8mm, minimum width=30mm] (rep) at (6.6,4) {historian replica};
  \node[hf/store, minimum height=8mm, minimum width=30mm] (his) at (6.6,3) {historian};
  \node[hf/reader, minimum height=7mm, minimum width=30mm] (hmi) at (6.6,2) {HMI / SCADA};
  \node[hf/source, minimum height=7mm, minimum width=30mm] (plc) at (6.6,1) {PLC / DCS};
  \node[hf/source, minimum height=7mm, minimum width=30mm] (fld) at (6.6,0) {sensors, valves};

  \draw[hf/flow, draw=hfSource] (fld) -- (plc);
  \draw[hf/flow, draw=hfSource] (plc) -- (hmi);
  \draw[hf/save] (hmi) -- (his);
  \draw[hf/save] (his) -- node[hf/note, right=4mm] {replicate} (rep);
  \draw[hf/ask] (off) -- (rep);

  \draw[hf/pain, dashed, -{Stealth[length=2mm]}] (9.6,4.75) -- (9.6,2.25);
  \node[text=hfAlert, font=\Large\bfseries] at (9.6,3.5) {$\times$};
  \node[hf/note, text=hfAlert, anchor=west] at (9.75,3.5) {no};
\end{tikzpicture}
\caption{A Purdue-style plant. Data flows up, the office reads a copy in the DMZ,
and nothing from level~4 reaches down into control.}
\end{figure}

\begin{concept}{The single reader}
One program reads the plant on everyone's behalf, keeps what it reads, and
becomes the \emph{source of truth} for everything above it. Everyone else asks
that program.

\textbf{What it fixes:} one connection to the plant, one copy of history, and
one program responsible for the data.

\textbf{What it costs:} that program is now on the live path. If it stops,
everyone goes blind.
\end{concept}

\section{Stop asking, start listening}

\begin{history}{1999 \textperiodcentered{} An oil pipeline and a satellite}
Andy Stanford-Clark at IBM and Arlen Nipper at Arcom had to bring data from oil
pipeline sensors into a SCADA system over satellite links. The links were billed
by the byte and often dropped. Polling over that was absurd, since each question
cost money and most answers were ``nothing changed''. They designed
\textbf{MQTT}.

It stayed niche for a decade. The open-source Mosquitto broker at the end of the
2000s and the rise of the ``Internet of Things'' changed that. OASIS approved
MQTT as a standard in 2014, and MQTT~5 followed in 2019.
\end{history}

\begin{concept}{Publish/subscribe}
A producer \emph{publishes} messages on a \emph{topic} to a \emph{broker}, only
when it has something to say. Consumers \emph{subscribe} to topics at the
broker. Producer and consumers never address each other.

\textbf{What it fixes:} the source publishes once, no matter how many readers
there are, and readers come and go without the source noticing.

\textbf{What it costs:} a broker to run, monitor and secure, and a message
format that every side has to agree on.
\end{concept}

\begin{figure}[htbp]
\centering
\begin{tikzpicture}
  \node[hf/source, minimum width=28mm] (dev) {pipeline sensor};
  \node[hf/bus, minimum width=24mm, minimum height=12mm, right=26mm of dev] (brk) {broker};
  \node[hf/reader, minimum width=26mm, right=26mm of brk, yshift=14mm] (r1) {SCADA};
  \node[hf/reader, minimum width=26mm, right=26mm of brk] (r2) {dashboard};
  \node[hf/store, minimum width=26mm, right=26mm of brk, yshift=-15mm] (r3) {archive};

  \draw[hf/push] (dev) -- node[hf/note, above] {publish once} node[hf/note, below] {on change} (brk);
  \draw[hf/push] (brk.east) -- (r1.west);
  \draw[hf/push] (brk.east) -- (r2.west);
  \draw[hf/push] (brk.east) -- (r3.west);
  \node[hf/note, above=1mm of brk, text=hfBus] {fan-out};
  \node[hf/note, below=1mm of brk] {topic\\\texttt{pipeline/7/pressure}};
\end{tikzpicture}
\caption{Publish/subscribe. The sensor doesn't know how many readers there are.}
\end{figure}

In industry, the idea has grown into the \emph{Unified Namespace}: one broker
where every system publishes its current state, and everyone else subscribes.
Other brokers carry the same idea, NATS being one, but the move is the same: a
bus in the middle decouples producers from consumers.

\section{Back to our two programs}

Now look again at the two programs from the start of the chapter. Each of the
three historical answers describes a different system.

\begin{figure}[htbp]
\centering
\begin{tikzpicture}[font=\sffamily\small]
  \tikzset{app/.style={hf/reader, minimum width=10mm, minimum height=7mm, inner sep=2pt}}
  % (a) everyone asks the source
  \begin{scope}
    \node[hf/source, minimum width=22mm] (pa) at (1.6,0) {plant};
    \foreach \i in {0,1,2} {
      \node[app] (a\i) at (0.4+1.2*\i,-2.6) {app};
      \draw[hf/ask] (a\i.north) -- (pa.south);
    }
    \node[hf/note, text width=40mm] at (1.6,-3.6) {\textbf{1.}\ everyone asks\\the source};
  \end{scope}
  % (b) single reader
  \begin{scope}[xshift=5.2cm]
    \node[hf/source, minimum width=22mm] (pb) at (1.6,0) {plant};
    \node[hf/store, minimum width=22mm, minimum height=8mm] (hb) at (1.6,-1.3) {historian};
    \draw[hf/ask] (hb.north) -- (pb.south);
    \foreach \i in {0,1,2} {
      \node[app] (b\i) at (0.4+1.2*\i,-2.6) {app};
      \draw[hf/ask] (b\i.north) -- (hb.south);
    }
    \node[hf/note, text width=40mm] at (1.6,-3.6) {\textbf{2.}\ one reader,\\everyone asks it};
  \end{scope}
  % (c) bus in the middle
  \begin{scope}[xshift=10.4cm]
    \node[hf/source, minimum width=22mm] (pc) at (1.6,0) {plant};
    \node[hf/bus, minimum width=22mm, minimum height=7mm] (bc) at (1.6,-1.3) {bus};
    \draw[hf/push] (pc.south) -- (bc.north);
    \foreach \i in {0,1,2} {
      \node[app] (c\i) at (0.4+1.2*\i,-2.6) {app};
      \draw[hf/push] (bc.south) -- (c\i.north);
    }
    \node[hf/note, text width=40mm] at (1.6,-3.6) {\textbf{3.}\ a bus\\in the middle};
  \end{scope}
\end{tikzpicture}
\caption{Three answers to ``who talks to the plant?''. Blue arrows ask, purple
arrows push.}
\end{figure}

\begin{enumerate}
  \item \textbf{Everyone asks the source.} This is the Modbus and early OPC
        world. It's simple and has no extra parts. The cost is load on the
        server that grows with every reader, a separate copy of the truth per
        reader, and a source that can't tell who depends on it.
  \item \textbf{One reader becomes the memory, and everyone asks it.} This is
        the historian and Purdue world. It has one connection to the plant, one
        copy of history and a clear owner. The cost is that the historian is now
        on the live path.
  \item \textbf{A bus in the middle.} This is the MQTT world. The source
        publishes once, any number of readers subscribe, and nobody depends on
        anybody's address. The cost is a broker to operate and a schema that
        all sides have to agree on.
\end{enumerate}

None of these is the modern one that made the others obsolete. All three are
still installed in plants today, often side by side. The question isn't which
one is newest. It's which pain you have now: load on the source, disagreeing
copies of history, or coupling between programs that should come and go freely.

Our two programs have a little of each.

\section*{Further reading}

\begin{itemize}
  \item J.~J.\ Downs and E.~F.\ Vogel, ``A plant-wide industrial process control
        problem'', \emph{Computers \& Chemical Engineering} 17(3), 245--255,
        1993.
  \item OPC Foundation, ``OPC Foundation History'' (2024),
        \url{https://opcfoundation.org/wp-content/uploads/2024/11/2024-11-12-OPCF-History.pdf}.
  \item E.~H.\ Bristol, U.S.\ Patent 4,669,097, ``Data compression for display
        and storage'', assigned to The Foxboro Company.
  \item HiveMQ, ``MQTT Essentials, Part 1: Introducing MQTT'', on the origins
        of MQTT.
\end{itemize}
